% v2 -- rewritten
% Seções:
%   6.1 - Modelagem do implante e tecido ósseo
%   6.2 - Simulação de cargas mastigatórias
%   6.3 - Validação K-Fold do modelo biomecânico
%   6.4 - Exercício prático: comparar 3 designs de implante
\destaque{Nivel-alvo N1 -- validacao de estresse em implantes usando Periomod e validacao cruzada K-Fold, dentro do Framework SUS-Twin.}
\label{ch:implantodontia}

\section{Modelagem do implante e tecido ósseo}
O planejamento digital de implantes dentários\footnote{Parafuso de titânio
inserido no osso maxilar para substituir a raiz de um dente perdido.} começa
com a CBCT, que gera um volume NIfTI. O fluxo SUS-Twin inicia com a
segmentação por densidade Hounsfield\footnote{Unidades Hounsfield (HU): água =
0~HU, osso cortical = 300--1000+~HU, ar = --1000~HU.} para separar osso
cortical e trabecular.

O código usa SimpleITK\footnote{Biblioteca Python de código aberto para
processamento de imagens médicas.} para ler a CBCT e extrair o sítio.

\begin{lstlisting}[language=Python,caption={Segmentação óssea com SimpleITK}]
import SimpleITK as sitk, numpy as np
img = sitk.ReadImage("cbct_paciente.nii")
arr = sitk.GetArrayFromImage(img)
cortical = np.where((arr >= 300) & (arr <= 1000), arr, 0)
trabecular = np.where((arr >= 100) & (arr < 300), arr, 0)
sitk.WriteImage(sitk.GetImageFromArray(cortical), "cortical.nii")
sitk.WriteImage(sitk.GetImageFromArray(trabecular), "trabecular.nii")

# Estatisticas de densidade ossea
print(f"Cortical: {np.mean(cortical[cortical>0]):.0f} HU")
print(f"Trabecular: {np.mean(trabecular[trabecular>0]):.0f} HU")
\end{lstlisting}

Esse mapeamento alimenta o modelo FEM no Periomod, permitindo que o SUS-Twin
atribua propriedades mecânicas distintas a cada região óssea e calcule a
distribuição de tensões iniciais.

\section{Simulação de cargas mastigatórias}
Definem-se dois cenários: (a) axial de 200~N simulando mastigação
molar\footnote{200~N equivalem a ~20~kg, força típica na mastigação.}; (b)
lateral de 50~N simulando bruxismo\footnote{Carga axial age ao longo do eixo
do implante; lateral age perpendicularmente, gerando momento fletor.}. O
Periomod resolve o campo de tensões por FEM.

\begin{lstlisting}[language=Python,caption={Condições de contorno para FEM}]
import numpy as np
axial = np.array([0.0, 0.0, -200.0])
lateral = np.array([50.0, 0.0, 0.0])
nos_oclu = np.arange(100, 200)
nos_base = np.arange(0, 50)
for no in nos_oclu: aplicar_forca(no, axial)
for no in nos_base: aplicar_restricao(no, [0, 0, 0])
print("Condicoes de contorno aplicadas: OK")
print(f"Forca axial: {np.linalg.norm(axial):.0f} N")
print(f"Forca lateral: {np.linalg.norm(lateral):.0f} N")
\end{lstlisting}

As tensões de von Mises são exportadas em MPa\footnote{MPa =
10\textsuperscript{6}~Pa; osso cortical fratura entre 100 e 200~MPa;
trabecular até ~10~MPa.} para a validação K-Fold, completando o ciclo
SUS-Twin de simulação biomecânica.

\section{Validação K-Fold do modelo biomecânico}
Para evitar overfitting, aplica-se K-Fold\footnote{K-Fold divide os dados em K
partes; treina-se em K--1 e testa-se na restante, repetindo K vezes.} com K=5
sobre os dados de CBCT. Um RandomForest\footnote{Conjunto de árvores de decisão
cuja predição final é a média de todas, reduzindo variância.} mapeia densidade
(HU) para tensão de pico (MPa). O SUS-Twin usa o R\textsuperscript{2} médio como métrica de
confiança para o plano cirúrgico.

\begin{lstlisting}[language=Python,caption={Validação K-Fold com RandomForest}]
from sklearn.model_selection import KFold
from sklearn.ensemble import RandomForestRegressor
import numpy as np

X = np.load("densidades_hu.npy")
y = np.load("tensoes_pico.npy")
kf = KFold(n_splits=5, shuffle=True, random_state=42)
modelo = RandomForestRegressor(n_estimators=100, max_depth=10)
scores = []

for fold, (tr, ts) in enumerate(kf.split(X), 1):
    modelo.fit(X[tr], y[tr])
    score = modelo.score(X[ts], y[ts])
    scores.append(score)
    print(f"Fold {fold}: R² = {score:.3f}")

print(f"R² medio: {np.mean(scores):.3f} ± {np.std(scores):.3f}")
\end{lstlisting}

O SUS-Twin persiste o modelo treinado e o associa ao plano cirúrgico,
permitindo que o clínico avalie probabilisticamente o risco de falha
mecânica antes da cirurgia.

\section{Exercício prático: comparar 3 designs de implante}
Execute a pipeline para três geometrias de implante\footnote{Cilíndrico: parede
reta. Cônico: afinamento apical. Rosqueado: filetes helicoidais que aumentam
o contato osso-implante.} e tabule os resultados.

\begin{table}[h]
\centering
\caption{Métricas comparativas entre designs de implante}
\label{tab:comparativo-implantes}
\begin{tabular}{lccc}
\toprule
\textbf{Design} & \textbf{Tensão (MPa)} & \textbf{Uniformidade (\%)} & \textbf{BIC (\%)}\\
\midrule
Cilíndrico & 142,3 & 62 & 38 \\
Cônico      & 118,7 & 78 & 45 \\
Rosqueado   & 97,2  & 84 & 52 \\
\bottomrule
\end{tabular}
\end{table}

Procedimento:
\begin{enumerate}
    \item Carregue cada STL no módulo de malha do Periomod.
    \item Atribua E: cortical 14~GPa, trabecular 1,5~GPa.
    \item Aplique carga axial de 200~N e execute a simulação FEM.
    \item Extraia a tensão de von Mises máxima e a distribuição.
    \item Calcule o BIC\footnote{Bone-implant contact (BIC): fração da superfície
    do implante em contato direto com osso; >40\% indica osseointegração.}
    a partir da malha deformada.
    \item Preencha a tabela e discuta o melhor compromisso no contexto
    do SUS-Twin.
\end{enumerate}

A comparação ilustra como a geometria impacta a longevidade do implante,
fechando o ciclo SUS-Twin de validação biomecânica baseada em evidências.
